%=====================================================================
\chapter{Sixteen-Week Teaching Schedule}\label{app:schedule}
%=====================================================================
\normalfigures

A mapping of the 22 chapters onto a sixteen-week semester of two lecture hours and
three laboratory hours a week --- the HEC 3 (2--3) pattern --- with assessment points.
Adjust to your own calendar; the dependencies are what matter.

\section{The Schedule}

\rowcolors{2}{white}{lightbg}
\begin{longtable}{@{}L{1.0cm}L{3.6cm}L{1.9cm}L{4.4cm}L{3.3cm}@{}}
\toprule
\rowcolor{primary!15}
\textbf{Wk} & \textbf{Lecture topic} & \textbf{Chapters} & \textbf{Lab} & \textbf{Assessment} \\
\midrule
\endhead
1 & What machine learning is; its mathematics & 1, 2 & A rule against a learned model;
the chapter in NumPy & --- \\
2 & The workflow and the toolkit & 3 & A leak-proof pipeline & --- \\
3 & Concept learning and the version space & 4 & Candidate elimination in forty lines &
Quiz 1 (Ch 1--3) \\
4 & Linear regression & 5 & Least squares three ways & --- \\
5 & Logistic regression & 6 & Logistic regression by hand and by library &
\textbf{Assignment 1 set} \\
6 & \textbf{Overfitting, bias--variance, regularisation} & 7 & Find the bottom of the U
& --- \\
7 & \textbf{Measuring classifier performance} & 8 & Evaluate an imbalanced classifier &
\textbf{Assignment 1 due}; Quiz 2 (Ch 4--7) \\
8 & Decision trees & 9 & ID3 by hand, pruning by cross-validation &
\textbf{Midterm (Ch 1--8)} \\
9 & Bayesian learning and naive Bayes & 10 & Naive Bayes by counting and by library &
--- \\
10 & Nearest neighbours; support vector machines & 11, 12 & Neighbours, scaling and $k$;
margins and kernels & \textbf{Assignment 2 set} \\
11 & Neural networks; the project brief & 13, 22 & Backpropagation from scratch &
\textbf{Project set}; proposal due \\
12 & Ensemble learning & 14 & One tree against five ensembles &
\textbf{Assignment 2 due} \\
13 & Clustering; PCA and self-organizing maps & 15, 16 & $k$-means and a dendrogram;
PCA and a SOM & Quiz 3 (Ch 9--14) \\
14 & Semi-supervised learning and EM; hidden Markov models & 17, 18 & EM by hand;
forward, Viterbi, Baum--Welch & Project experiments frozen \\
15 & Markov decision processes; reinforcement learning & 19, 20 & Value and policy
iteration; bandits and Q-learning & --- \\
16 & Responsible machine learning; project presentations & 21 & Auditing a model for
fairness & \textbf{Project report and presentation} \\
\bottomrule
\end{longtable}

\begin{alertbox}[If You Must Cut Something]
Cut \emph{depth}, never Chapters 7 and 8. Generalisation and evaluation are the hinge
of the course: a student who cannot tell a good model from a lucky one will misreport
every later result. If time is short, compress Chapter 4 to the version space and
candidate elimination worked example, teach Chapter 16's self-organizing map in one
lecture, and combine Chapters 19 and 20 --- but keep the Q-learning update, which the
HEC outline's reinforcement-learning outcome requires.
\end{alertbox}

\section{Dependency Map}

\dsfigh{%
\begin{tikzpicture}[font=\scriptsize,
  n/.style={rounded corners=2pt, draw=#1, fill=#1!10, text=#1!50!black,
            minimum width=1.05cm, minimum height=0.5cm, font=\tiny\bfseries},
  crit/.style={rounded corners=2pt, draw=accent, line width=1.4pt, fill=accent!16,
               text=accent!70!black, minimum width=1.05cm, minimum height=0.5cm,
               font=\tiny\bfseries},
  a/.style={-{Stealth[length=1.6mm]}, gray!55, line width=0.6pt}]
% Part I
\foreach \i/\c in {1/1,2/2,3/3,4/4}{
  \pgfmathsetmacro{\xx}{(\i-1)*1.35}
  \ifnum\c=3 \node[crit] (c\c) at (\xx,2.4) {\c};
  \else \node[n=primary] (c\c) at (\xx,2.4) {\c}; \fi}
\node[font=\tiny, anchor=west, text=primary] at (8.4,2.4) {Part I --- foundations};
% Part II
\foreach \i/\c in {1/5,2/6,3/7,4/8}{
  \pgfmathsetmacro{\xx}{(\i-1)*1.35}
  \ifnum\c>6 \node[crit] (c\c) at (\xx,1.55) {\c};
  \else \node[n=secondary] (c\c) at (\xx,1.55) {\c}; \fi}
\node[font=\tiny, anchor=west, text=secondary] at (8.4,1.55) {Part II --- linear models, generalisation};
% Part III
\foreach \i/\c in {1/9,2/10,3/11,4/12,5/13,6/14}{
  \pgfmathsetmacro{\xx}{(\i-1)*1.35}
  \node[n=greenhl] (c\c) at (\xx,0.7) {\c};}
\node[font=\tiny, anchor=west, text=greenhl!50!black] at (8.4,0.7) {Part III --- supervised algorithms};
% Part IV
\foreach \i/\c in {1/15,2/16,3/17}{
  \pgfmathsetmacro{\xx}{(\i-1)*1.35}
  \node[n=purplehl] (c\c) at (\xx,-0.15) {\c};}
\node[font=\tiny, anchor=west, text=purplehl] at (8.4,-0.15) {Part IV --- without full labels};
% Part V
\foreach \i/\c in {1/18,2/19,3/20}{
  \pgfmathsetmacro{\xx}{(\i-1)*1.35}
  \node[n=tealhl] (c\c) at (\xx,-1.0) {\c};}
\node[font=\tiny, anchor=west, text=tealhl] at (8.4,-1.0) {Part V --- sequences and decisions};
% Part VI
\foreach \i/\c in {1/21,2/22}{
  \pgfmathsetmacro{\xx}{(\i-1)*1.35}
  \node[n=goldhl] (c\c) at (\xx,-1.85) {\c};}
\node[font=\tiny, anchor=west, text=goldhl!50!black] at (8.4,-1.85) {Part VI --- practice};
% key dependencies
\draw[a] (c2) -- (c5);
\draw[a] (c3) -- (c7);
\draw[a] (c4) -- (c9);
\draw[a] (c6) -- (c13);
\draw[a] (c9) -- (c14);
\draw[a] (c10) -- (c17);
\draw[a] (c15) -- (c17);
\draw[a] (c17) -- (c18);
\draw[a] (c18) -- (c19);
\draw[a] (c19) -- (c20);
\draw[a] (c8) .. controls +(4.2,0) and +(0,1.2) .. (c21);
\draw[a] (c7) .. controls +(5.2,-0.3) and +(0.3,1.6) .. (c22);
\end{tikzpicture}}%
{Chapter dependencies. The outlined chapters (3, 7 and 8) are the ones every later
chapter relies on; the arrows mark the strongest direct dependencies.}{dependency-map}

\section{Laboratory Sessions}

Each week's three laboratory hours follow the same pattern: the chapter's
\emph{Try It Yourself} lab first (about one hour, from the runnable file in
\texttt{code/}), then the chapter's applied review questions worked in pairs, then ---
from Week~11 --- project time with the instructor circulating. Every lab runs on a
standard Anaconda installation (Python 3.10 or later, NumPy, pandas, SciPy,
scikit-learn); none needs a GPU or an internet connection.
